\section{Conclusion}\label{sec:conclusion}
Part-of-speech tags and lemmas were put into a subword tokenizer in several forms and measured under one protocol in five languages, with gold morpheme boundaries, held-out text, repeated training samples and a language-model probe.
Weighting BPE merges by part-of-speech transitions, the method the authors first proposed, changes neither the boundaries a tokenizer finds nor the text it produces, because a transition probability describes the order of words and not their inside, and weighting merges inside content words changes the vocabulary in the intended direction with an effect on accuracy that depends on the language.
Deriving boundaries from lemmas raises boundary recall substantially everywhere at a cost of four to eight percent more tokens, with a gain on unseen words that depends on how distinctive the affixes of the language are, so a tokenizer that wants morphological boundaries should obtain them from a lemmatizer or a segmentation lexicon, and a tagger adds nothing on top.
What those boundaries are worth to a language model is a separate matter, since the probe ranks the tokenizers in a different order from the boundary scores, with the unigram model first and Lemma2Word last in every language, and the two orders will have to be reconciled at a larger scale than a 7M-parameter model before morphological alignment can be recommended as a design goal.
